\chapter{Conclusão}
\label{cap:conclusao}

Este trabalho apresentou o desenvolvimento e a avaliação da VANSI, uma ferramenta interativa projetada para auxiliar o processo de ensino de algoritmos de análise sintática. A motivação para este projeto surgiu da complexidade inerente à disciplina de Compiladores e da carência de ferramentas que ofereçam uma visualização clara e passo a passo dos algoritmos de análise sintática.

\section{Considerações finais}

Os objetivos propostos inicialmente foram alcançados. A ferramenta foi desenvolvida utilizando tecnologias modernas de desenvolvimento \textit{web} e suporta os algoritmos \gls{ll}(1), \gls{slr} e \gls{lr}(1). A avaliação da ferramenta foi realizada com um grupo de alunos da Universidade Federal do Ceará.

A análise quantitativa dos resultados demonstrou uma alta taxa de satisfação e intenção de uso futuro. As avaliações positivas entre a avaliação geral e a intenção de recomendação indica que a VANSI possui um valor percebido significativo pelos usuários. Além disso, a análise dos padrões de interação revelou que os alunos exploraram as diferentes funcionalidades de visualização.

\section{Limitações}

Apesar dos objetivos propostos terem sido alcançados, é importante reconhecer as limitações deste trabalho. A amostra de participantes foi relativamente pequena e composta por alunos voluntários. Além disso, a avaliação foi realizada sem grupo de controle, o que pode limitar a generalização dos resultados. A avaliação focou na satisfação e na intenção de uso, sem uma análise aprofundada do impacto da ferramenta no desempenho acadêmico dos alunos. Por fim, a avaliação não considerou a opinião de docentes sobre a ferramenta, o que poderia fornecer \textit{insights} para melhorias futuras.

\section{Trabalhos futuros}

Considerando os resultados obtidos e as limitações deste trabalho, ainda há oportunidades de expansão para este estudo. Como trabalhos futuros, sugere-se a realização de uma avaliação mais profunda com grupos de controle, a inclusão de suporte para analisadores mais complexos, a implementação de uma \textit{engine} de animação utilizando \gls{svg} para melhorar a performance das animações e, por fim, a otimização da interface para uso em \textit{smartphones}.